\documentclass{article}
\usepackage[utf8]{inputenc}
\usepackage{hyperref}
\usepackage{url}
\title{?????/QSPR ????? transport ????????(Tang?? ?????????? ??????)?? ???????? ???????? ?????? ????????? Tang}
\author{sci-adk (deterministic render)}
\date{\today}

\begin{document}
\maketitle

\noindent\textit{Draft compiled by sci-adk from Spec \texttt{pfas-rice-oos-tang} (v1). Belief state is revisable as Evidence accrues.}

\section{Goal}
?????/QSPR ????? transport ????????(Tang?? ?????????? ??????)?? ???????? ???????? ?????? ????????? Tang
2026?? ???????? ???????????? ????????? ?????????? (?????????? ??????? log10 RMSE?? in-sample Tang-????????
??????? ??????????, ??????? ???? ?????????).

\section{Background}
????? sci-adk run??? "Yamazaki ????? = ???????????"??? REFUTED?? ???????????(???? in-sample ??????;
???????? RMSE 0.84). ?????? ????? Yamazaki ???? ??????? ????? ??????????. ????? ?????????? ?????????
???????? **??????? ???? ?????? ?????? ?????????**?????? ???????(out-of-sample) ?????????. Tang 2026???
Yamazaki?? ????? ?????(????? paddy pot)$\cdot$??????(Nipponbare)$\cdot$??????? ?????????? ????????? ???????? ?????????
(TF, stalk/leaf/endosperm)??? ???????????, ?????/QSPR(monotone) transport ??????????? Tang??
?????????? ????? ?????? ??????? ????? ????? ??????? ???????? ?????.

\section{Method}
`model\_api.tang\_tf\_validation`(ORYZA2000 biomass, dry-weight TF)?? PFOA$\cdot$PFOS$\cdot$GenX 3?????
???? f\_xy\_source="recommended"(????? monotone, Tang ????????)?? ????? TF??? Tang ?????? TF?? ????????,
?????? Tang-???????? f\_xy(in-sample)?? ???????????(`validation/oos\_tang.py`, 0.1 $\mu$g/g ????????).
??????? log10 RMSE?? in-sample RMSE??? ?????????. ??????(Tang 2026)?? ????? ???????.

Planned approaches:
\begin{itemize}
  \item `model\_api.tang\_tf\_validation`(ORYZA2000 biomass, dry-weight TF)?? PFOA$\cdot$PFOS$\cdot$GenX 3?????
???? f\_xy\_source="recommended"(????? monotone, Tang ????????)?? ????? TF??? Tang ?????? TF?? ????????,
?????? Tang-???????? f\_xy(in-sample)?? ???????????(`validation/oos\_tang.py`, 0.1 $\mu$g/g ????????)
  \item ??????? log10 RMSE?? in-sample RMSE??? ?????????
  \item ??????(Tang 2026)?? ????? ???????.
\end{itemize}

\section{Hypotheses and findings}

\subsection{?????/QSPR ????? transport ????????(Tang?? ?????????? ??????)?? ???????? ???????? ?????? ????????? Tang
2026?? ???????? ???????????? ????????? ?????????? (?????????? ??????? log10 RMSE?? in-sample Tang-????????
??????? ??????????, ??????? ???? ?????????).}
\begin{itemize}
  \item Hypothesis id: \texttt{hyp-001} (exploratory)
  \item Decision rule (qualitative): Expert judgment based on evidence
  \item \textbf{Status: refuted} --- confidence strong (graded)
  \item Basis: qualitative judge: Out-of-sample prediction of the independent Tang 2026 dataset FAILS: log10 RMSE 1.23 (\textasciitilde{}17x) with theory params vs 0.52 only after an in-sample Tang refit. The miss is systematic (PFSA under, ether over), so f\_xy is dataset/condition-dependent and the model does not predict one dataset from another's calibration. Hypothesis refuted -- consistent with hyp-yamazaki (in-sample reproduction != out-of-sample prediction).
  \item \textit{Evidence validity:} referent=empirical; data\_source(s)=measured
\end{itemize}

\section{Evidence}
\begin{itemize}
  \item \texttt{evi-oos-tang} (experiment\_run): point=1.232, finding=Out-of-sample vs the INDEPENDENT Tang 2026 dataset (theory/QSPR monotone f\_xy, NOT fit to Tang): log10 RMSE 1.232 (\textasciitilde{}17x)
\end{itemize}

\section{Pending agent judgments}
These hypotheses use a non-numeric rule and await an in-session agent verdict (no autonomous LLM call):
\begin{itemize}
  \item \texttt{hyp-001} (qualitative): Expert judgment based on evidence
  \item finding: Out-of-sample vs the INDEPENDENT Tang 2026 dataset (theory/QSPR monotone f\_xy, NOT fit to Tang): log10 RMSE 1.232 (\textasciitilde{}17x). Systematic miss -- PFSA way UNDER (PFOS stalk 0.013 vs 0.571, endosperm 0.004 vs 0.944; \textasciitilde{}40-200x) and ether OVER (GenX stalk 8.3 vs 0.92, leaf 17.7 vs 1.33; \textasciitilde{}10-13x). Only an IN-SAMPLE Tang-refit f\_xy reaches 0.519. The model does NOT predict an independent dataset from another dataset's parameters: f\_xy is dataset/condition-dependent (PFSA head-group offset, GenX ether offset).
\end{itemize}

\section{References}
No literature cited.

\end{document}